\oldpage[0\textsubscript{IV-1}]{111}

\begin{definition}[19.8.4]
\label{0.19.8.4-es}
Sean $A$ un anillo local, $P\to A$ el único homomorfismo de un anillo local
primo $P$ en $A$, y $p$ la característica de los cuerpos residuales de $P$ y
$A$. Se dice que $A$ es un \emph{anillo de Cohen} si es una $P$-álgebra de
Cohen, es decir (19.8.1), si:
\begin{enumerate}
  \item[$1^{\mathrm o}$] $A$ es noetheriano y completo.
  \item[$2^{\mathrm o}$] $A$ es un $P$-módulo plano (lo que equivale también a
  decir que $A$ es un $\widehat P$-módulo plano (Bourbaki, \emph{Alg. comm.},
  cap.~III, §~5, n\textsuperscript{o}~4, prop.~4)).
  \item[$3^{\mathrm o}$] $A/pA$ es un cuerpo (necesariamente separable sobre el
  cuerpo residual de $P$, pues este cuerpo es primo).
\end{enumerate}
\end{definition}

Si $p=0$, estas condiciones equivalen a decir que $A$ es un \emph{cuerpo de
característica $0$}. Si $p>0$, se tiene necesariamente $pA\neq0$; la condición
$3^{\mathrm o}$ significa que $pA$ es el ideal maximal $\mathfrak m$ de $A$;
la condición $2^{\mathrm o}$ significa que $p$ es $A$-regular, puesto que $P$
es un anillo de valoración discreta ($\mathbf 0_{\mathrm I}$, 6.3.4). Por tanto,
$A$ es un anillo regular (17.1.1, $d)$) de dimensión $1$ y, en consecuencia, un
anillo de valoración discreta, completo en virtud de $1^{\mathrm o}$; en suma:

\begin{proposition}[19.8.5]
\label{0.19.8.5-es}
Los anillos de Cohen son los cuerpos de característica $0$ y los anillos de
valoración discreta completos cuyo cuerpo residual tiene característica
$p>0$ y cuyo ideal maximal está generado por $p\cdot1$ (siendo $1$ la unidad del
anillo).
\end{proposition}

Obsérvese que, en el segundo caso, $p\cdot1\neq0$ porque $p$ es $A$-regular;
por tanto, se puede identificar $p\cdot1$ con el entero $p$; el homomorfismo
canónico $\mathbf Z_p\to A$ es inyectivo, y se identifica $p\cdot1$ con el
elemento $p$ de $\mathbf Z_p$; en este caso se dice que $A$ es un
\emph{$p$-anillo de Cohen}.

\begin{theorem}[19.8.6]
\label{0.19.8.6-es}
\textup{(Cohen).} ---
\begin{enumerate}
  \item[(i)] Sean $W$ un anillo de Cohen, $C$ un anillo local noetheriano
  completo y $\mathfrak J$ un ideal de $C$ distinto de $C$. Entonces todo
  homomorfismo local $u:W\to C/\mathfrak J$ se factoriza como
  $W\xrightarrow{v}C\to C/\mathfrak J$, donde $v$ es un homomorfismo local.
  \item[(ii)] Sea $K$ un cuerpo. Existe un anillo de Cohen $W$ cuyo cuerpo
  residual es isomorfo a $K$. Si $W'$ es un segundo anillo de Cohen y $K'$ su
  cuerpo residual, todo isomorfismo $u:K\xrightarrow{\sim}K'$ proviene, por
  paso a los cocientes, de un isomorfismo $v:W\xrightarrow{\sim}W'$.
\end{enumerate}
\end{theorem}

\begin{proof}
Esto no es sino (19.8.2) aplicado al caso en que $A$ es un anillo local primo.
\end{proof}

\begin{remarks}[19.8.7]
\label{0.19.8.7-es}
\begin{enumerate}
  \item[(i)] Cuando $K$ es de característica $0$, la parte (ii) de (19.8.6)
  se vuelve trivial.
  \item[(ii)] El homomorfismo $v$ de (19.8.6, (i)) \emph{no está
  necesariamente determinado de manera única por $u$}, como muestra ya el
  caso en que $W$ es un cuerpo de característica $0$, $\mathfrak J^2=0$ y $u$
  es un isomorfismo (cf.~(21.5.5)). Del mismo modo, en (19.8.6, (ii)), el
  isomorfismo $v$ \emph{no está necesariamente determinado de manera única
  por $u$} (cf.~(21.5.5)).

  Sin embargo, cuando $K$ es \emph{perfecto} y de característica $p>0$, se
  verá en (21.5.5) que en (19.8.6, (ii)) el isomorfismo $v$ es único. Se verá
  también más adelante que, en este caso, $W$ se identifica con el anillo
  $W_\infty(K)$ de \emph{vectores de Witt de longitud infinita} sobre $K$.
  \item[(iii)] En (19.8.6, (i)), se pueden debilitar las hipótesis sobre $C$
  utilizando (19.3.10) y (19.3.12).
\end{enumerate}
\end{remarks}

\begin{theorem}[19.8.8]
\label{0.19.8.8-es}
\textup{(Cohen).} --- Sean $A$ un anillo local noetheriano completo y $k$ su
cuerpo residual.

\oldpage[0\textsubscript{IV-1}]{112}

\begin{enumerate}
  \item[(i)] Existe un anillo de Cohen $W$ tal que $A$ es isomorfo a un anillo
  cociente de un anillo de series formales $W[[T_1,\ldots,T_n]]$ (y, en
  particular, $A$ es isomorfo a un cociente de un anillo local regular
  completo (17.3.8)). Si $A$ contiene un cuerpo, es isomorfo a un anillo
  cociente de $k[[T_1,\ldots,T_n]]$.
  \item[(ii)] Supongamos además que $A$ es íntegro. Entonces existe un
  subanillo $B$ de $A$ tal que: $1^{\mathrm o}$ $B$ es isomorfo a un anillo de
  series formales sobre un anillo $C$ que es un cuerpo o un anillo de Cohen
  (lo que implica que $B$ es un anillo local regular y completo (17.3.8));
  $2^{\mathrm o}$ $B$ tiene el mismo cuerpo residual que $A$ y la inyección
  $B\to A$ es un homomorfismo local; $3^{\mathrm o}$ $A$ es una $B$-álgebra
  finita.
\end{enumerate}
\end{theorem}

\begin{proof}
Sea $\mathfrak m$ el ideal maximal de $A$. Existe un anillo de Cohen $W$ cuyo
cuerpo residual es isomorfo a $k$ (19.8.6, (ii)); se tiene, pues, un
homomorfismo local $W\to A/\mathfrak m$, que se factoriza entonces como
$W\xrightarrow{u}A\to A/\mathfrak m$, donde $u$ es un homomorfismo local
(19.8.6, (i)). Para toda familia finita $(x_i)_{1\leq i\leq n}$ de elementos
de $\mathfrak m$, existe entonces un homomorfismo local
$v:W[[T_1,\ldots,T_n]]\to A$ que prolonga $u$ y tal que $v(T_i)=x_i$ para todo
$i$ (Bourbaki, \emph{Alg. comm.}, cap.~III, §~4, n\textsuperscript{o}~5,
prop.~6). Cuando $A$ contiene un cuerpo, contiene un cuerpo primo $P$, del cual
$k$ es una extensión (necesariamente separable) y, por tanto, $A$ contiene un
cuerpo isomorfo a $k$ (19.6.2); se puede entonces sustituir $W$ por $k$ en la
definición anterior de $v$.

\medskip
\noindent(i)\enspace
Tomemos primero como $x_i$ un sistema de generadores de $\mathfrak m$. Como
$W$ tiene el mismo cuerpo residual que $A$, y como las clases de los $x_i$ en
el anillo graduado $\gr_\bullet(A)$ generan $\gr_\bullet(A)$ como $k$-álgebra,
$\gr(v):\gr_\bullet(W[[T_1,\ldots,T_n]])\to\gr_\bullet(A)$ es sobreyectivo;
se deduce que $v$ es también sobreyectivo (Bourbaki, \emph{Alg. comm.},
cap.~III, §~2, n\textsuperscript{o}~8, cor.~2 del teorema~1). Recordemos que el
caso en que $A$ contiene un cuerpo ya se ha tratado y solo figura aquí como
recordatorio (19.6.3).

\medskip
\noindent(ii)\enspace
Si $A$ contiene un cuerpo, contiene un cuerpo $k'$ isomorfo a $k$, como se ha
visto; se considera entonces un sistema de parámetros
$(y_j)_{1\leq j\leq m}$ de $A$ (16.3.6), se toma
$B=k[[T_1,\ldots,T_m]]$ y se considera el homomorfismo local $w:B\to A$ cuya
restricción a $k$ es un isomorfismo $k\to k'$ y tal que
$w(T_j)=y_j$ para $1\leq j\leq m$. Si $A$ no contiene ningún cuerpo, el único
homomorfismo $\mathbf Z_{p\mathbf Z}\to A$ (19.8.3) es necesariamente
inyectivo (de lo contrario, como $A$ es íntegro, su núcleo sería el ideal
maximal de $\mathbf Z_{p\mathbf Z}$ y su imagen sería isomorfa a un cuerpo);
además, se tiene entonces $p>0$ por hipótesis, pues
$\mathbf Z_{p\mathbf Z}$ es un cuerpo si $p=0$. El elemento $p\cdot1$ de $A$
(identificado con $p$) no es divisor de cero en $A$ y pertenece a
$\mathfrak m$; por tanto, (16.3.4 y 16.3.7), existe una familia
$(z_i)_{1\leq i\leq m-1}$ que, junto con $p$, forma un sistema de parámetros
de $A$. El anillo de Cohen $W$ considerado al comienzo de la demostración es
entonces un anillo de valoración discreta de cuerpo residual $k$, en el que
$p$ genera el ideal maximal (19.8.5), y el homomorfismo único $u:W\to A$
definido al comienzo envía $p$ a sí mismo. Se toma entonces
$B=W[[T_1,\ldots,T_{m-1}]]$ y se considera el homomorfismo local $w:B\to A$
que coincide con $u$ al restringirlo a $W$ y tal que $w(T_i)=z_i$ para
$1\leq i\leq m-1$. En ambos casos, si $\mathfrak n$ es el ideal maximal de
$B$, está claro que $\mathfrak nA$ es un ideal de definición de $A$; como,
además, $B/\mathfrak n$ y $A/\mathfrak m$ son isomorfos, $A$ es un
$B$-módulo cuasifinito ($\mathbf 0_{\mathrm I}$, 7.4.4), y, por tanto, un
$A$-módulo de tipo finito, pues $B$ es completo y $A$ es
separado para las topologías $\mathfrak n$-preádicas
($\mathbf 0_{\mathrm I}$, 7.4.1). Por otra
parte,

\oldpage[0\textsubscript{IV-1}]{113}

en ambos casos se tiene $\dim(B)=\dim(A)=m$; en el primer caso, esto resulta
de (17.1.4, (iii)); en el segundo, se ve directamente que $p$ y los $T_i$
($1\leq i\leq m-1$) forman una sucesión $B$-regular que genera $\mathfrak n$,
o también se puede utilizar el hecho de que estos elementos generan
$\mathfrak n$ y que se tiene $\dim(B)\geq\dim(A)$ por (16.3.10). Como $A$ y
$B$ son íntegros, se deduce finalmente de (16.3.10) que $w$ es inyectivo, lo
que termina la demostración.
\end{proof}

\begin{corollary}[19.8.9]
\label{0.19.8.9-es}
Sea $A$ un anillo local noetheriano íntegro completo que contiene un cuerpo
$k_0$; sea $k$ el cuerpo residual de $A$, y supongamos que $k$ es finito sobre
$k_0$. Entonces, en la conclusión de (19.8.8, (ii)), se pueden sustituir
$1^{\mathrm o}$ y $2^{\mathrm o}$ por la condición de que $B$ sea de la forma
$k_0[[T_1,\ldots,T_n]]$, siendo la inyección canónica $B\to A$ un
$k_0$-homomorfismo local (para la estructura usual de $k_0$-álgebra de $B$).
\end{corollary}

\begin{proof}
En efecto, retomando la demostración de (19.8.8, (ii)), se define esta vez
$w:k_0[[T_1,\ldots,T_n]]\to A$ de modo que su restricción a $k_0$ sea la
identidad y envíe $T_j$ a $y_j$ para $1\leq j\leq n$. La hipótesis de que $k$ es de
grado finito sobre $k_0$ implica de nuevo que $A$ es un $B$-módulo cuasifinito,
y, por tanto, de tipo finito por ($\mathbf 0_{\mathrm I}$, 7.4.1), y se
concluye como en (19.8.8).
\end{proof}

\begin{corollary}[19.8.10]
\label{0.19.8.10-es}
Sea $A$ un anillo local artiniano cuyo ideal maximal $\mathfrak m$ es de
cuadrado nulo; existe entonces un anillo local noetheriano regular $B$, de
ideal maximal $\mathfrak n$, tal que $A$ es isomorfo a
$B/\mathfrak n^2$.
\end{corollary}

\begin{proof}
Sean $K$ el cuerpo residual $A/\mathfrak m$ de $A$ y $n$ el rango de
$\mathfrak m/\mathfrak m^2=\mathfrak m$ sobre $K$. Si $A$ contiene un cuerpo,
se deduce de (19.6.3) que $A$ es isomorfo a $B/\mathfrak b$, donde
$B=K[[T_1,\ldots,T_n]]$ y $\mathfrak b$ está contenido en el cuadrado del
ideal maximal $\mathfrak n$ de $B$; pero, como
$\operatorname{long}(B/\mathfrak n^2)=n+1=\operatorname{long}(A)$, se tiene
necesariamente $\mathfrak b=\mathfrak n^2$.

Supongamos ahora que $A$ no contiene ningún cuerpo; esto implica que $K$ es
de característica $p>0$ y que $p\cdot1\neq0$ en $A$ (19.6.3); por tanto,
$p\cdot1$ es un elemento de $\mathfrak m$, y existen otros $n-1$ elementos
$x_i$ ($2\leq i\leq n$) de $\mathfrak m$ que, junto con $p\cdot1$, forman una
base de $\mathfrak m$ sobre $K$. Sea $W$ un anillo de Cohen cuyo cuerpo
residual es isomorfo a $K$; $W$ es un anillo de valoración discreta cuyo ideal
maximal está generado por $p$; se ha visto en la demostración de (19.8.8) que
existe un homomorfismo $u:W\to A$ que envía $p$ a sí mismo y que, por paso
a los cocientes, da la identidad sobre $K$. Se toma
$B=W[[T_2,\ldots,T_n]]$ y se considera el homomorfismo local $w:B\to A$ que
coincide con $u$ al restringirlo a $W$ y tal que $w(T_i)=x_i$ para $2\leq i\leq n$. Está
claro que $w$ es sobreyectivo y que su núcleo $\mathfrak b$ está contenido en
el cuadrado del ideal maximal $\mathfrak n=pB+BT_2+\cdots+BT_n$ de $B$; como
$\operatorname{long}(B/\mathfrak n^2)=n+1=\operatorname{long}(A)$, se tiene
de nuevo $\mathfrak b=\mathfrak n^2$.
\end{proof}

\begin{proposition}[19.8.11]
\label{0.19.8.11-es}
Sean $A$ un anillo local artiniano, $\mathfrak m$ su ideal maximal y $k$ su
cuerpo residual. Para que $A$ sea isomorfo a un anillo cociente de un anillo
de Cohen, es necesario y suficiente que $\mathfrak m$ esté generado por
$p\cdot1$, donde $p$ es la característica de $k$.
\end{proposition}

\begin{proof}
La condición es evidentemente necesaria (19.8.5). Para ver que es suficiente,
se observa, como al comienzo de la demostración de (19.8.8), que existe un
anillo de Cohen $W$ cuyo cuerpo residual es isomorfo a $k$ y un homomorfismo
local $u:W\to A$. Además, si se considera el homomorfismo compuesto
$\mathbf Z_{p\mathbf Z}\to W\xrightarrow{u}A$ (que es necesariamente el único
homomorfismo de $\mathbf Z_{p\mathbf Z}$ en $A$), se ve que la imagen

\oldpage[0\textsubscript{IV-1}]{114}

por $u$ del elemento $p\cdot1$ de $W$ es el elemento $p\cdot1$ de $A$; como
el elemento $p\cdot1$ de $W$ genera el ideal maximal de este anillo, se deduce
inmediatamente de la hipótesis que
$\gr(u):\gr_\bullet(W)\to\gr_\bullet(A)$ es sobreyectivo y, por tanto, también
lo es $u$ (Bourbaki, \emph{Alg. comm.}, cap.~III, §~2,
n\textsuperscript{o}~8, cor.~2 del teorema~1).
\end{proof}
\subsection{Álgebras relativamente formalmente lisas}
\label{subsection:0.19.9-es}

\begin{definition}[19.9.1]
\label{0.19.9.1-es}
Sean $\Lambda$ un anillo topológico, $A$ una $\Lambda$-álgebra topológica y
$B$ una $A$-álgebra topológica. Se dice que $B$ es una $A$-álgebra
\emph{formalmente lisa relativamente a $\Lambda$} si, para toda $A$-álgebra
topológica discreta $C$ y todo ideal nilpotente $\mathfrak J$ de $C$, todo
$A$-homomorfismo continuo $u_0:B\to C/\mathfrak J$ que se factoriza como
$B\xrightarrow{u}C\to C/\mathfrak J$, donde $u$ es un $\Lambda$-homomorfismo
continuo, se factoriza también como $B\xrightarrow{v}C\to C/\mathfrak J$, donde
$v$ es un $A$-homomorfismo continuo.
\end{definition}

De esta definición resulta que, si $B$ es una $A$-álgebra \emph{formalmente
lisa}, entonces $B$ es también \emph{formalmente lisa relativamente a
$\Lambda$}, para toda estructura de $\Lambda$-álgebra topológica definida sobre
$A$ (es decir, todo homomorfismo continuo de anillos $\Lambda\to A$).

\begin{proposition}[19.9.2]
\label{0.19.9.2-es}
Sean $\Lambda$ un anillo topológico y $A$ una $\Lambda$-álgebra topológica.
\begin{enumerate}
  \item[(i)] $A$ es una $\Lambda$-álgebra formalmente lisa relativamente a
  $\Lambda$.
  \item[(ii)] Si $B$ es una $A$-álgebra formalmente lisa relativamente a
  $\Lambda$ y $C$ una $B$-álgebra formalmente lisa relativamente a $\Lambda$,
  entonces $C$ es una $A$-álgebra formalmente lisa relativamente a $\Lambda$.
  \item[(iii)] Sean $B$ una $A$-álgebra formalmente lisa relativamente a
  $\Lambda$ y $A'$ una $A$-álgebra topológica; entonces la $A'$-álgebra
  topológica $B\otimes_A A'$ es formalmente lisa relativamente a $\Lambda$.
  \item[(iv)] Sean $B$ una $A$-álgebra topológica y $S$ (resp. $T$) un
  subconjunto multiplicativo de $A$ (resp. $B$) tal que la imagen canónica de
  $S$ en $B$ esté contenida en $T$. Si $B$ es una $A$-álgebra formalmente lisa
  relativamente a $\Lambda$, entonces $T^{-1}B$ es una $S^{-1}A$-álgebra
  formalmente lisa relativamente a $\Lambda$.
  \item[(v)] Sean $B_i$ ($1\leq i\leq n$) unas $A$-álgebras topológicas. Para
  que $\prod_{i=1}^n B_i$ sea una $A$-álgebra formalmente lisa relativamente a
  $\Lambda$, es necesario y suficiente que cada una de las $B_i$ lo sea.
\end{enumerate}
\end{proposition}

\begin{proof}
La afirmación (i) es trivial, y la demostración de las demás sigue muy de cerca
las demostraciones de (19.3.5); se deja, pues, al lector.
\end{proof}

\begin{corollary}[19.9.3]
\label{0.19.9.3-es}
Sean $A$ un anillo topológico, y $A'$ y $B$ dos $A$-álgebras topológicas.
Entonces la $A$-álgebra topológica $B'=B\otimes_A A'$ es una $B$-álgebra
formalmente lisa relativamente a $A$.
\end{corollary}

\begin{proof}
Esto resulta de (19.9.2, (i) y (iii)).
\end{proof}

\begin{proposition}[19.9.4]
\label{0.19.9.4-es}
Sean $\Lambda$ un anillo topológico, $A$ una $\Lambda$-álgebra topológica y
$B$ una $A$-álgebra topológica. Las condiciones siguientes son equivalentes:
\begin{enumerate}
  \item[a)] $B$ es una $A$-álgebra formalmente lisa relativamente a $\Lambda$.
  \item[b)] $\widehat B$ es una $A$-álgebra formalmente lisa relativamente a
  $\Lambda$.
  \item[c)] $\widehat B$ es una $\widehat A$-álgebra formalmente lisa
  relativamente a $\Lambda$.
  \item[d)] $\widehat B$ es una $\widehat A$-álgebra formalmente lisa
  relativamente a $\widehat\Lambda$.
\end{enumerate}
\end{proposition}

\begin{proof}
Se deja de nuevo al lector la demostración, calcada de la de (19.3.6).
\end{proof}

\oldpage[0\textsubscript{IV-1}]{115}

\begin{env}[19.9.5]
\label{0.19.9.5-es}
Del mismo modo, el enunciado (19.3.8) sigue siendo válido (con la misma
demostración) cuando se sustituye «formalmente lisa» por «formalmente lisa
relativamente a $\Lambda$». Si en el enunciado de (19.3.10) se sustituye
«formalmente lisa» por «formalmente lisa relativamente a $\Lambda$», la
conclusión se sustituye por la siguiente (permaneciendo la demostración
esencialmente inalterada): \emph{todo $A$-homomorfismo continuo
$u_0:B\to C/\mathfrak J$ que se factoriza como
$B\xrightarrow{u}C\to C/\mathfrak J$, donde $u$ es un $\Lambda$-homomorfismo
continuo, se factoriza también como $B\xrightarrow{v}C\to C/\mathfrak J$, donde
$v$ es un $A$-homomorfismo continuo}.
\end{env}

\begin{env}[19.9.6]
\label{0.19.9.6-es}
Los criterios de lisura formal (19.4.1) y (19.4.2) son válidos cuando se
sustituye en ellos «formalmente lisa» por «formalmente lisa relativamente a
$\Lambda$», permaneciendo las demostraciones prácticamente inalteradas.
\end{env}

\begin{proposition}[19.9.7]
\label{0.19.9.7-es}
Sean $\Lambda$ un anillo topológico, $A$ una $\Lambda$-álgebra topológica y
$B$ una $A$-álgebra topológica. Supongamos que, para toda $A$-álgebra discreta
$C$ y todo ideal $\mathfrak J$ de $C$ tal que $\mathfrak J^2=0$, todo
$A$-homomorfismo continuo $u_0:B\to C/\mathfrak J$ que se factoriza como
$B\xrightarrow{u}C\to C/\mathfrak J$, donde $u$ es un $\Lambda$-homomorfismo
continuo, se factoriza también como $B\xrightarrow{v}C\to C/\mathfrak J$, donde
$v$ es un $A$-homomorfismo continuo. Entonces $B$ es una $A$-álgebra
formalmente lisa relativamente a $\Lambda$.
\end{proposition}

\begin{proof}
La demostración de (19.4.3) se transcribe inmediatamente.
\end{proof}

\begin{proposition}[19.9.8]
\label{0.19.9.8-es}
Sean $\Lambda$ un anillo topológico, $A$ una $\Lambda$-álgebra topológica y
$B$ una $A$-álgebra topológica. Para que $B$ sea una $A$-álgebra formalmente
lisa relativamente a $\Lambda$, es necesario y suficiente que, para todo
$B$-módulo topológico discreto $L$, anulado por un ideal abierto de $B$, se
tenga (cf.~18.4.2)
\[
  \operatorname{Exalcotop}_{A/\Lambda}(B,L)=0.
\]
\end{proposition}

\begin{proof}
Con las notaciones de la demostración de (19.4.4), basta observar aquí que se
puede suponer que la extensión $E_\Lambda$ de $B_\Lambda$ es
$\Lambda$-trivial; el resto de la demostración permanece entonces inalterado.
\end{proof}

Cuando $\Lambda$, $A$ y $B$ son anillos discretos, el criterio (19.9.8) se
reduce a
\begin{equation}
\tag{19.9.8.1}
\label{0.19.9.8.1-es}
  \operatorname{Exalcom}_{A/\Lambda}(B,L)=0
  \quad\text{para todo $B$-módulo $L$}~;
\end{equation}
es decir, \emph{toda $A$-extensión conmutativa de $B$ por un $B$-módulo, que es
$\Lambda$-trivial, es también $A$-trivial}.

\subsection{Álgebras formalmente no ramificadas y álgebras formalmente
étales}
\label{subsection:0.19.10-es}
\label{0.19.10-es}

\begin{env}[19.10.1]
\label{0.19.10.1-es}
También tendremos que utilizar dos nociones próximas a la noción de álgebra
formalmente lisa.
\end{env}

\begin{definition}[19.10.2]
\label{0.19.10.2-es}
Sean $A$ un anillo topológico y $B$ una $A$-álgebra topológica. Se dice que
$B$ es una $A$-álgebra \emph{formalmente no ramificada} si, para todo
$A$-homomorfismo $p:E\to C$ de $A$-álgebras discretas cuyo núcleo es
nilpotente, y todo $A$-homomorfismo continuo $u:B\to C$, existe \emph{a lo sumo
un} $A$-homomorfismo continuo $v:B\to E$ tal que $u$ se factoriza como
$B\xrightarrow{v}E\xrightarrow{p}C$. Se dice que $B$ es una $A$-álgebra
\emph{formalmente étale} si es formalmente lisa y formalmente no ramificada, es
decir, si, para todo $A$-homomorfismo sobreyectivo $p:E\to C$ de $A$-álgebras
discretas cuyo núcleo es nilpotente, y para todo $A$-homomorfismo continuo
$u:B\to C$, existe \emph{uno y solo un} $A$-homomorfismo continuo $v:B\to E$
tal que $u$ se factoriza como $B\xrightarrow{v}E\xrightarrow{p}C$.
\end{definition}

\oldpage[0\textsubscript{IV-1}]{116}

\begin{examples}[19.10.3]
\label{0.19.10.3-es}
Nos limitaremos a los anillos \emph{discretos}.
\begin{enumerate}
  \item[(i)] Si el homomorfismo estructural $\rho:A\to B$ es
  \emph{sobreyectivo}, $B$ es una $A$-álgebra \emph{formalmente no ramificada},
  pues, con las notaciones de (19.10.2), el homomorfismo compuesto
  $A\xrightarrow{j}B\xrightarrow{v}E$ debe ser el homomorfismo estructural, por
  lo que $v$ es único (si existe). En cambio, $B$ no es formalmente lisa salvo
  que $\rho$ sea \emph{biyectivo}.
  \item[(ii)] Sean $A$ un anillo y $S$ un subconjunto multiplicativo de $A$;
  entonces $B=S^{-1}A$ es una $A$-álgebra \emph{formalmente étale}. En efecto,
  sean $j:A\to B$ el homomorfismo canónico, $\rho:A\to E$ un homomorfismo de
  anillos, $\mathfrak J$ un ideal nilpotente de $E$ y
  $p:E\to C=E/\mathfrak J$ el homomorfismo canónico. Si $u:B\to C$ es un
  $A$-homomorfismo, $u(s/1)$ es invertible en $C$ para $s\in S$, puesto que
  $s/1$ es invertible en $B$; pero $u(s/1)=p(\rho(s))$, y, como $\mathfrak J$
  es nilpotente, el hecho de que la imagen de $\rho(s)$ en $E/\mathfrak J$ sea
  invertible implica que el propio $\rho(s)$ es invertible en $E$
  ($\mathbf 0_{\mathrm I}$, 7.1.12); se concluye que $\rho$ se factoriza de una
  sola manera como $A\xrightarrow{j}B\xrightarrow{v}E$, y se tiene
  evidentemente $u\circ j=(p\circ v)\circ j$, de donde $u=p\circ v$, lo que
  prueba nuestra afirmación.
  \item[(iii)] Una extensión finita separable $k'$ de un cuerpo $k$ es una
  $k$-álgebra formalmente lisa (19.6.1); veremos más adelante que, para que una
  extensión de \emph{tipo finito} $K$ de un cuerpo $k$ sea una $k$-álgebra
  formalmente étale, es necesario y suficiente que sea formalmente no
  ramificada, y esto equivale también a que $K$ sea una extensión finita y
  separable de $k$ (21.7.4).
  \item[(iv)] Un álgebra de polinomios $A[X_\alpha]_{\alpha\in I}$ sobre un
  anillo $A$ es \emph{formalmente lisa} sobre $A$ (19.3.3), pero \emph{no
  formalmente étale} si $I$ no es vacío, como resulta inmediatamente de las
  definiciones.
\end{enumerate}
\end{examples}

\begin{remark}[19.10.4]
\label{0.19.10.4-es}
El mismo razonamiento que en (19.4.3) muestra que, para verificar que una
$A$-álgebra topológica $B$ es formalmente no ramificada, se puede, en la
definición (19.10.2), limitar al caso en que el núcleo de $p:E\to C$ es de
\emph{cuadrado nulo}; se puede además limitar al caso en que $u(B)\subset p(E)$
(pues de lo contrario no existe ningún $v$ que factorice $u$) y, sustituyendo
$C$ por $u(B)$ y $E$ por $p^{-1}(u(B))$, suponer $u$ y $p$ sobreyectivos.
\end{remark}
